\honest{Knowing About Biases Can Hurt People}{Eliezer Yudkowsky}{2007}

I once told my mother that when an expert says they are 99\% confident, they are right only about 70\% of the time. Then I remembered who I was talking to, and added that she must apply this skepticism to herself too. She said, ``Are you kidding? This is great! I'm going to use it all the time!''

I give no source for the numbers. \nb{The best-known study with those numbers, by Fischhoff, Slovic and Lichtenstein, asked students general-knowledge questions. They were not experts, and some experts, such as weather forecasters, are well calibrated. So the post's first sentence hands the mother a weapon stronger than the evidence for it.}

Then one study: Taber and Lodge, on students' views about affirmative action and gun control. It lists six predictions and calls all of them confirmed. People rate arguments for what they already believe more highly, even when told to be objective. They spend more effort attacking contrary arguments than supporting ones. They seek out sources that agree with them. A balanced set of arguments for and against makes them more polarized. People with stronger attitudes show all this more. And people who know more about politics show it more, because they have more ammunition for counterarguing. \nb{Taber and Lodge themselves note that attempts to reproduce the polarization have often failed. The sixth finding, the one that carries the title, is about knowing more about politics. It says nothing about knowing about biases.}

``If you're irrational to start with, having more knowledge can hurt you.'' For an ideal Bayesian, information never has negative expected value; humans can cut themselves.

I have seen people ``severely messed up'' by their knowledge of biases: it gives them more ammunition against anything they dislike. Too much ready ammunition is one of the main ways mentally agile people end up stupid, in Keith Stanovich's sense of ``dysrationalia.'' Before telling you about any of them, I ask you, ``You can think of people who fit this description, right?''

Then two stories about one man. He learned that experts are often overconfident, concluded that you cannot trust experts, and went off into his own ``highly questionable extrapolation.'' The biases he had learned seemed much less salient when they applied to his own conclusions. That is the lesson I gave my mother. Then I told him about sophisticated arguers, and the next time I said something he disliked, he called me one without naming any flaw. He ``had acquired yet another Fully General Counterargument.''

So in my last talk I introduced biases with the conjunction fallacy and then spent thirty minutes hammering on confirmation bias, motivated skepticism and sophisticated arguers. A simple example is enough to interest an audience. But ``The literature on bias is mostly cognitive psychology for cognitive psychology's sake,'' so if I did not give the warnings in that one lecture, they would probably never hear them.

My rule is to never mention calibration and overconfidence before talking about disconfirmation bias and sophisticated arguers. ``First, do no harm!'' \nb{The rule changes only the order of the warnings. The man in the story had heard the warning about sophisticated arguers, and used it against the author. The post itself says that notion ``can be deadly.''}
